\section{Related Work}

\subsection{Application Caches and Concurrent Updates}

Nishtala et al.\ describe Memcache leases used to prevent stale sets when a client refills a key after a concurrent update \cite{nishtala2013}. A lease and a per-item generation both give a reader a value to present when it attempts to refill a key. The POS implementation advances an invalidation generation for each recipe that depends on a changed ingredient. It also compares stock versions for fills within the same generation. Our focus is the application logic that connects an order transaction, a dependent item, and a later display.

Atikoglu et al.\ observed substantial differences in access patterns among application cache pools \cite{atikoglu2012}. In a POS workload, skewed item selection can favor cache reuse, while shared ingredients expand the set of keys removed after a sale. Cao and Liu reported similar average client response times for adaptive TTL and invalidation in their web accelerator. They also reported a higher worst-case delay in their invalidation implementation because the accelerator blocked new requests while sending invalidations over TCP \cite{cao1998}. That observation motivates measuring order acknowledgments separately from reads. It describes the implementation in their study. Franklin et al.\ similarly show that locality and conflicts affect the cost of supporting transactional client caches \cite{franklin1997}.

\subsection{Consistency and Measurement}

TxCache provides a consistent but potentially older database snapshot for transactions that use cached values \cite{ports2010}. CacheGenie uses database triggers to maintain cached application queries \cite{gupta2011}. FlightTracker provides read-your-writes behavior across read-optimized stores with different replication and caching arrangements \cite{shi2020}. In the POS stress test, the largest order version recorded before a GET provides a lower bound for that GET. A later response that meets that bound can still miss a more recent commit. In the controlled stockout test, all browsing starts after stock is known to be zero, so the comparison has a direct business interpretation.

Harizopoulos et al.\ identify transaction-processing overheads that remain even when a database fits in memory \cite{harizopoulos2008}. Li et al.\ discuss the effect of scheduling and shared resources on tail latency \cite{li2014}. Our measurements count SQL calls alongside GET and POST latency so that reductions in menu queries can be reported independently of latency differences. The Yahoo! Cloud Serving Benchmark separates workload mix and record selection from system policy \cite{cooper2010}. Schroeder et al.\ show that independently timed arrivals can expose different behavior from completion-driven clients \cite{schroeder2006}. This motivates the use of open-loop k6 arrivals for response time and delivery. The separate correctness exercises use specified events and concurrency because their purpose is to observe availability transitions.
